% !TEX program = xelatex
% 编译方式：xelatex main.tex （连续编译两次以生成目录；建议用 latexmk -xelatex）
% 依赖宏包：ctex / xeCJK / tcolorbox / enumitem / booktabs / array / longtable / xcolor / fancyhdr / hyperref
%
% 【填写指引】搜索 TODO 字样，替换封面与页眉中的标题、作者、出处即可。
% 章节写在 chapters/ 下，ch00 起逐章 \input。

\documentclass[UTF8,a4paper,12pt,fontset=windows]{ctexart}

% ---------- 页面与行距 ----------
\usepackage[a4paper,left=2.6cm,right=2.6cm,top=3.0cm,bottom=3.2cm]{geometry}
\renewcommand{\baselinestretch}{1.32}
\setlength{\parskip}{0.2em}

% ---------- 中文字体 ----------
% Windows 自带字体；macOS 可改为 Songti SC / PingFang SC，Linux 装 Noto CJK 后用 fontset=none 另设
\setCJKfamilyfont{yahei}{Microsoft YaHei}
\newcommand{\yh}{\CJKfamily{yahei}}

% ---------- 常用宏包 ----------
\usepackage{amsmath,amssymb}
\usepackage{enumitem}
\usepackage{booktabs}
\usepackage{array}
\usepackage{longtable}
\setlength{\tabcolsep}{4pt}
\emergencystretch=3em
\raggedbottom
\usepackage{xcolor}
\usepackage{tcolorbox}
\tcbuselibrary{breakable,skins}
\usepackage{fancyhdr}
\usepackage[colorlinks=true,linkcolor=teal,urlcolor=teal]{hyperref}

% ---------- 配色（可整体换色） ----------
\definecolor{ink}{RGB}{40,44,52}
\definecolor{teal}{RGB}{15,84,96}
\definecolor{clay}{RGB}{186,98,74}
\definecolor{mist}{RGB}{243,246,244}

% ---------- 章节标题 ----------
\ctexset{
  part = {
    format = {\centering\yh\huge\bfseries\color{teal}},
    aftertitle = {\par\vspace{4pt}{\color{clay}\rule{0.4\linewidth}{1.4pt}}\par},
    beforeskip = {3em},
    afterskip = {2em},
  },
  section = {
    format = {\yh\Large\bfseries\color{teal}},
    beforeskip = {2.4em},
    afterskip = {0.8em},
  },
  subsection = {
    format = {\yh\normalsize\bfseries\color{ink}},
    beforeskip = {1.6em},
    afterskip = {0.5em},
  },
  subsubsection = {
    format = {\yh\normalsize\bfseries\color{ink}},
    beforeskip = {1.2em},
    afterskip = {0.4em},
  },
}

% ---------- 自定义组件 ----------
% 核心洞察框
\newtcolorbox{insight}{
  breakable, enhanced,
  colback=mist, colframe=teal, boxrule=0.8pt,
  arc=1.5mm, left=2.2mm, right=2.2mm, top=1.6mm, bottom=1.6mm,
  before skip=6pt, after skip=8pt
}
% 提醒 / 边界框
\newtcolorbox{notehot}{
  breakable, enhanced,
  colback=white, colframe=clay, boxrule=0.8pt,
  arc=1.5mm, left=2.2mm, right=2.2mm, top=1.6mm, bottom=1.6mm,
  before skip=6pt, after skip=8pt
}
% 清单环境（□ 复选框）
\newenvironment{checklist}
  {\begin{itemize}[leftmargin=2em,itemsep=1.5pt,topsep=3pt,label={$\square$}]}
  {\end{itemize}}
% 术语强调
\newcommand{\term}[1]{\textbf{\color{teal}#1}}
% 核心一句
\newcommand{\key}[1]{\par\vspace{2pt}\noindent\textcolor{teal}{\textbf{核心一句：}}#1\par\vspace{4pt}}

% ---------- 页眉页脚 ----------
\pagestyle{fancy}
\fancyhf{}
\setlength{\headheight}{16pt}
\fancyhead[L]{\small\yh\color{ink}讲义标题} % TODO
\fancyhead[R]{\small\yh\color{ink}\leftmark}
\fancyfoot[C]{\small\color{ink}\thepage}
\renewcommand{\headrulewidth}{0.3pt}

\hypersetup{pdftitle={讲义标题},pdfauthor={作者或整理者}} % TODO

\begin{document}

% ---------- 封面（TODO 全部替换） ----------
\begin{titlepage}
  \centering
  \vspace*{2.2cm}
  {\yh\Huge\bfseries\color{teal} 讲义主标题\par}
  \vspace{0.5em}
  {\yh\LARGE\bfseries\color{ink} 讲义副标题\par}
  \vspace{1.8em}
  {\large\color{ink} 一条贯穿全文的分析轴（可选）\par}
  {\large\color{ink} 关键词 · 关键词 · 关键词\par}
  \vspace{2.8cm}
  {\color{ink} 思想内核提炼自某平台博主「博主名」系列 N 期视频\par}
  {\color{ink} 在其观点与经验之上，整理延伸为一份完整讲义\par}
  \vfill
  {\large 2026 年 \quad 整理版\par}
  \vspace{1.5cm}
  {\color{clay} 写给目标读者的一句话\par}
\end{titlepage}

% ---------- 目录 ----------
\pagestyle{empty}
\pagenumbering{roman}
\tableofcontents
\clearpage
\pagestyle{fancy}
\pagenumbering{arabic}

% ---------- 正文（章节文件放在 chapters/ 下） ----------
\section*{写在前面}
\addcontentsline{toc}{section}{写在前面}

\subsection*{这份讲义是什么}

这是一份\term{观点讲义}：把某位博主在一系列视频里反复表达的世界观、概念与方法，提炼、重组、延伸成一份可以从头读到尾的连贯文档。它不是视频字幕的摘抄合集，而是一次再创作——素材属于博主，结构属于整理者。

\subsection*{这份讲义不是什么}

\begin{checklist}
  \item 不是学术综述，也不是专业建议。涉及医学、法律、投资等领域时，博主的经验只能作为视角，不能替代专业人士；
  \item 不是逐字稿。口语被改写成书面表达，例子被重新编排，但观点不替博主发明；
  \item 不是「正确答案」。凡是博主本人标注为猜测、个体经验的内容，讲义同样保留这份不确定。
\end{checklist}

\subsection*{怎么读}

\key{先接受一套新的解释框架，再决定拿走哪些方法。} 第一部分通常回答「为什么会这样」（认知层），中间部分回答「那怎么办」（方法层），最后是延伸思考与清单。可以顺序读，也可以直接跳到与你处境最相关的一章。

\section{第一章：核心判断与底层模型}

% 写作提示（正式成稿后删除本注释）：
% 本章回答「这件事在底层是怎么回事」。先用一个反直觉判断或场景开场，
% 再展开博主的 1~2 个底层模型/隐喻，每个模型标注出处 BV，
% 全章安排 1~2 个 insight 框，章末用 \key{} 收束并把问题交给下一章。

\subsection{从一个熟悉的场景开始}

在这里放一个目标读者一看就「对，我就是这样」的场景，然后给出博主的判断。

\begin{insight}
用 \term{博主原词} 表达的核心模型：一句话写清这个模型如何解释上面的场景。
\end{insight}

\subsection{模型如何运转}

展开机制：它在什么条件下成立、和「通常以为」的解释差别在哪。
关键论断后标注出处，如（BVxxxx）。整理者自己的延伸推理用
「顺着这个逻辑……」之类的语气带出，与博主原话区分开。

\key{章末金句：把这一章压缩成读者能转述给别人的一句话。}

% \input{chapters/ch02}
% …按实际章节增减…
\appendix

\section*{附录 A：核心概念速查表}

\begin{center}
\small
\begin{tabular}{>{\raggedright\arraybackslash}p{3.2cm} p{12.0cm}}
\toprule
\textbf{概念} & \textbf{含义} \\
\midrule
概念一 & 用一句话写清博主的定义，保留他本人的用词 \\
概念二 & …… \\
\bottomrule
\end{tabular}
\end{center}

\clearpage
\section*{附录 B：出处对照}

本讲义的思想内核提炼自博主的 N 期视频。下表为主要观点与原视频的对照，方便回看原作：

\begin{small}
\begin{longtable}{>{\raggedright\arraybackslash}p{5.2cm} >{\raggedright\arraybackslash}p{10.0cm}}
\toprule
\textbf{视频（标题）} & \textbf{本讲义中的相关章节 / 概念} \\
\midrule
\endfirsthead
\toprule
\textbf{视频（标题）} & \textbf{相关章节 / 概念（续）} \\
\midrule
\endhead
\bottomrule
\endfoot
示例标题 & 第 X 章某节、某概念 \\
\end{longtable}
\end{small}

\vspace{1em}

\term{观点属于博主，编排属于整理者；如有曲解，责任在整理一方。} 请把这份讲义当作一张地图，而不是一份结论——地图的作用是让你看清地形，路终究要自己走。


\end{document}
